% Extracted from attraction_wfigs.docx; see PLAN.md and tools/extract_pilot.py.

% DOCX paragraph 55
\chapter{Gravity … Attracts!}\label{ch:01}

% DOCX paragraph 58
As I have already discussed, the title of this chapter is meant to be less trivial than
it might seem.

% DOCX paragraph 59
Here, I do not want to simply stress the existence of a physical ``force'' exerted
between two objects with a mass, which attracts them to each other even if they are at a
great distance. Indeed, in chapter 3, we will discover that this idea, however
widespread and easy to understand, is incorrect and, at least in part, misleading.
Rather, with this title, I wish to emphasise that there is something --- gravity ---
that attracts us in a broad sense, and above all, attracts our attention. In addition to
physical objects, it applies an irresistible grip on our imagination; it can push our
thoughts towards radically new horizons, different from the usual ones, and expose us to
scenarios that extend far beyond our everyday experience.

% DOCX paragraph 60
But let's proceed step by step. To better define what gravity is, it helps to break down
the knowledge we have into three distinct but interconnected levels. In particular, we
could say that we have: \emph{instinctive knowledge}\index{knowledge!instinctive},
\emph{rational knowledge}\index{knowledge!rational}\emph{,} and finally, an
\emph{imaginative knowledge}\index{knowledge!imaginative} of gravity.

% DOCX paragraph 61
Let's see what they are and how they differ from each other.

% DOCX paragraph 64
\section{Instinctive Knowledge}\label{sec:01-instinctive}

% DOCX paragraph 66
We all know that instinct is the natural and intrinsic tendency (that is to say, not
requiring the intervention of reasoning or reflection) to carry out a particular
behaviour. An example of instinct would be what drives us to flinch to protect our heads
when we are surprised by a sudden, loud, and unknown sound. When put in these terms, it
is difficult to believe that something instinctive or irrational can possibly bind us to
gravity. Yet it is so.

% DOCX paragraph 67
Anyone who has handled a newborn will surely have seen first-hand the so-called
Moro\index{Moro!reflex} reflex, named after the Austrian paediatrician Ernst
Moro\index{Moro!Ernst} (1874-1951). It is one of the main neonatal reflexes and is
widely used in evaluating the function of the central nervous system. To induce this
reflex, it is sufficient to take a newborn, even in the first few seconds after birth,
lift it in a horizontal position, and (making sure it is totally safe) create the
sensation of a free fall\index{free fall}! The newborn's response to the unexpected and
apparent loss of support is a reaction of surprise, accompanied by the sudden opening of
the arms and hands in search of a handhold, as is shown in Figure \ref{fig:01-01}.

% DOCX paragraph 68
From a medical point of view, the appearance of this reflex is important evidence of the
perfect physiological functioning of the infant's central nervous system. That is why we
have all been subjected to this test. Any parent who has attended --- this has happened
to me three times --- is well aware of the relief they felt on seeing their son or
daughter respond as they should to this rather bizarre stimulus.

% DOCX paragraph 70
\begin{figure}[htbp]
\centering
\reflectbox{\includegraphics[width=0.65\linewidth]{../figs/source/en/fig-01-01.jpg}}

% DOCX paragraph 71
\caption{\emph{Example of Moro reflex: a newborn reacts instinctively to the loss of
support, spreading her arms and trying to grab a handhold to stop her fall. © V.
Tverdokhlib/depositphotos.com.}}\label{fig:01-01}
\end{figure}

% DOCX paragraph 73
On an anthropological level, Moro's reflex reminds us of our past as primates, when in
all probability we had to be immediately ready to follow our mothers, holding on firmly
to their backs. However, what interests us most here is its physical significance. The
presence of this instinct\index{knowledge!instinctive} a few seconds after birth (when
we are completely defenceless and know nothing about the world around us) reveals an
important truth concerning our interaction with gravity: we know it instinctively, well
before having conscious interactions with the rest of the physical universe. Thus, after
having spent nine comfortable months in our mothers' wombs, essentially isolated from
everything, we are immediately able to respond to gravity (or rather, to its absence).
This is no small detail.

% DOCX paragraph 74
The Moro reflex disappears at around six months, so our knowledge of gravity, while
remaining partly instinctive, changes over time as we develop the ability to observe the
physical universe and understand its laws.

% DOCX paragraph 77
\section{Rational Knowledge}\label{sec:01-rational}

% DOCX paragraph 79
As our experience of the world expands and our intellectual abilities are refined, our
knowledge of gravity changes from an instinctive level to a rational one. In other
words, it becomes an integral part of our expectations of how the world around us
functions. This is clearly demonstrated through simple visual experiments conducted on
very young children using cartoons. Although some of the subjects involved are not yet
able to walk, all of them can interpret the motion of an object that is consistent with
the presence of a gravitational field. The classic example is that of a sphere rolling
on a table: children react differently (at the level of facial expressions and eye
movements) depending on whether the sphere, once it reaches the edge of the table,
falls, or continues its motion despite being without support or, even more bizarrely, if
it begins to fly. The fact that the infants know what the correct behaviour of the ball
should be is further confirmation of how deeply the knowledge of the ``force'' of
gravity is rooted in our minds.

% DOCX paragraph 80
This plays a fundamental role in our rational perception\index{perception} of reality.
Thanks to this profound influence, our brains are able to solve very complex dynamic
problems in a very short time and with practically no effort. \index{knowledge!rational}
A simple example is the action of hurrying down a flight of stairs: it is one of the
most complex problems when programming robots (and one at which they often fail),
however, we humans can cope without even consciously thinking about it. Yet,
establishing in what sequence and at what speed our movements must be performed to
ensure the subtle balance between the various forces competing with the gravitational
force, is by no means trivial.

% DOCX paragraph 81
Finally, another property of gravity is worth pondering: its ability to stimulate our
imagination.

% DOCX paragraph 84
\section{Imaginative Knowledge}\label{sec:01-imaginative}

% DOCX paragraph 86
If it is clear that we have an instinctive and rational notion of gravity, then in my
opinion, it is equally clear that gravity also exerts an irresistible attraction on our
imagination. It is precisely because we are immersed and subjected to a gravitational
field every instant of our lives that we are naturally fascinated by those scenarios in
which gravity is weak or even absent. Who has never wished they could jump off a high
cliff or a mountain top and \ldots{} fly? Who has never imagined being an astronaut
aboard the International Space Station or a character in a science fiction movie
floating effortlessly from one place to another in zero gravity? I often do ... In other
words, gravity attracts our attention. It stimulates our imagination precisely because
it is the only ``force'' of which we have a conscious awareness and because we know all
too well how difficult it is to escape it. What, if not imagination, led first Isaac
Newton\index{Newton!Isaac} and then Albert Einstein\index{Einstein!Albert} to explain
(in very different ways) the laws that govern gravity?

% DOCX paragraph 87
There are many examples that illustrate the powerful pull of gravity on our imagination,
but I will limit myself to just one, which I find representative and easy to understand.
In 2013, the Spanish director Alfonso Cuarón\index{Cuarón, Alfonso (""Gravity"")} made a
movie with the emblematic title: \emph{Gravity}. For almost two hours, the film talks
about nothing other than gravity; or rather, its absence. However, not many know that
\emph{Gravity} broke the box office record for the first weekend screening of the films
released that autumn. Some may argue that this success is entirely due to the leading
actors, the two Hollywood superstars Sandra Bullock and George Clooney. In my opinion,
however, a fundamental role has been played by the fact that --- like it or not --- we
cannot escape gravity and its irresistible pull on our minds.
\index{knowledge!imaginative}

% DOCX paragraph 90
\section{One of Four, but Definitely Different from the Others}\label{sec:01-interactions}

% DOCX paragraph 92
The considerations made so far provide an excellent starting point for introducing
another important reflection on gravity's role within our understanding of nature.
Modern physics teaches us that there are four types of fundamental interactions, which
essentially describe all the processes taking place in the universe: the
\emph{electromagnetic interaction}\index{interaction!electromagnetic}, the \emph{strong
interaction}\index{interaction!strong}, the \emph{weak
interaction}\index{interaction!weak}, and the
\emph{gravitational}\index{interaction!gravitational} \emph{interaction}.

% DOCX paragraph 93
The first, the electromagnetic interaction, among other things, allows you to read this
book regardless of the format you are using. In fact, electromagnetic
waves\index{waves!electromagnetic} (photons\index{photon}, or more simply, light)
emanate from the page in front of you and reach your eyes. These waves then convert into
electrical signals transmitted via the optic nerve to the brain, which (thanks to a
complex combination of electrical and chemical exchanges) translates them into the words
you are reading. The electromagnetic interaction is also responsible for the cohesion
and dynamics of the molecules that compose us. Without it, we would not even exist as
human beings, and our molecules would scatter like bits of paper in the wind. The theory
that describes the electromagnetic interaction is well known, both at the level of
classical physics (where it is expressed by the Maxwell equations) and at the level of
quantum physics when dealing with elementary particles (where it is also known as theory
of \emph{quantum-electrodynamics}).

% DOCX paragraph 94
The second, the strong interaction, develops instead on the smallest scale accessible to
us in nature, of the order of a few fermi (or femtometres), that is, a few thousandths
of a billionth of a millimetre. The strong interaction is a hundred times more intense
than the electromagnetic interaction and is exerted between \emph{quarks}, the
constituent parts of elementary particles such as protons or neutrons. In reality, this
interaction is also present on slightly larger scales, i.e., inside
atomic\index{atomic!nucleus} nuclei (which generally have dimensions of the order of 10
fermi), where it takes the name of the \emph{strong nuclear force}. In the case of
protons and neutrons, the mediating particles of the strong interaction are called
\emph{gluons}\index{gluons}, while in the strong nuclear force, they are represented by
the \emph{pions}\index{pions}. In essence, the strong interaction acts as the ``glue''
that holds the atomic nuclei together, from the smallest atoms (hydrogen) to the larger
ones (for example, uranium). It also regulates what happens when two protons smash
against each other at a speed almost equal to that of light or when a neutron
star\index{neutron star} is born (a phenomenon we will discuss in detail in Chapter 5).
The theory that describes this interaction is well developed and is called
\emph{quantum-chromodynamics}\index{quantum-chromodynamics}. Unfortunately, given the
complexity of the theory and the equations that describe it, it is often difficult to
make precise predictions, especially if the energies are high or if the number of
particles involved is large, as in the case of neutron stars.

% DOCX paragraph 95
The third, the weak interaction, is responsible for the radioactive decay of some atomic
nuclei and acts between \emph{leptons}\index{leptons} (a class of particles to which
electron\index{electron}s, certainly the most ``familiar'' of the group, belong) and
quarks. It is thanks to this interaction that neutrinos\index{neutrino} (very light
particles produced by matter of high density and temperature, such as those at the
centre of the Sun\index{Sun}) interact only rarely (``weakly'') with protons and
neutrons, the ordinary matter of which we are composed. It is good to recall that at
this precise moment, our bodies are crossed by billions of neutrinos emitted about eight
minutes ago by the Sun and which travelled to us at almost the speed of light. If we do
not ``feel'' pierced by these particles, it is precisely because they interact only
weakly with the ordinary (or \emph{hadronic}) matter\index{hadronic (ordinary) matter}
of which we are made. In other words, there is little to worry about; we are essentially
transparent to neutrinos. The theory that describes the weak interaction is also very
well developed and can be combined with that of electromagnetism, so that we really talk
of the \emph{electroweak interaction}.

% DOCX paragraph 96
So, finally, we arrive at the fourth, the gravitational interaction. Although it is not
yet the time to explain in detail what it consists of and how it links to one of the
most subtle and elegant concepts of theoretical physics --- the curvature of spacetime
--- we can already reflect on what distinguishes it from the others. Gravity is the only
physical interaction of which we are consciously aware. The only one whose influence on
your body you are aware of at this moment, as you read this book. In fact, whether you
are lying on a bed, slumped in an armchair, or standing, \emph{you know} that some
``force'' affects your position. Without it, you would levitate, floating freely just
like an astronaut on the International Space Station.

% DOCX paragraph 97
This point is essential in order to fully understand the title of the chapter and grasp
its deeper meaning: gravity attracts our attention because we can directly and tangibly
experience its existence, unlike other fundamental interactions. By contrast, we cannot
appreciate the level of cohesion of the molecules that make up our bodies, how rarely we
interact with a neutrino, or how many radioactive particles we
emit.\textsuperscript{\hyperlink{note:01-01}{1}}

% DOCX paragraph 98
The fact that we can experience it directly and consciously is enough to give gravity a
special place among the fundamental interactions and, in my opinion, places it above all
others. What makes it even more peculiar is this awareness accompanies us in every
second of our lives, even before coming into the world. Even subconsciously, we know
that gravity exists long before grazing our knees when we learn to walk and run, and
certainly before we encounter the laws of physics in school or college.

% DOCX paragraph 101
\section{But What is Gravity, and How Does It Work?}\label{sec:01-gravity}

% DOCX paragraph 103
In all likelihood, many of you feel convinced you can provide a reasonable answer to the
simple questions: ``What is gravity, and how does it work?''. This is because you know
it both instinctively and rationally, and (either at school or university) you have been
given an ``educated'' explanation for it. However, it is equally likely that the
explanations given to you are not entirely correct, even if they are not completely
far-fetched either. What I mean is that what you have been told or taught isn't really
\emph{wrong}, but it's not \emph{right} either.

% DOCX paragraph 104
The reason for this apparent contradiction is that it is possible to understand gravity
on various levels. For example, the representation of gravity proposed by Newton
\index{Newton!Isaac}is simpler and more intuitive, which involves the existence of a
\emph{gravitational}\index{gravitational!force} \emph{force} and whose mathematical
description is relatively simple. At the same time, a different and deeper understanding
is possible: the one suggested by Einstein, which implies a geometric vision of space
and time and a much more complex mathematical description.

% DOCX paragraph 105
In the following chapters, a sort of evolution awaits us; we will develop our
understanding of gravity. The first stage of this evolution --- and therefore the first
destination of our virtual journey --- will take us from the basic understanding of
gravity as it is encoded in our brains, even on an instinctive level, to the description
provided by Newton's theory\index{Newton!Isaac}. The second stage --- or the next
destination of the journey --- will instead lead us to the mathematically compact and
physically profound description Einstein proposed with his theory of general
relativity\index{relativity!general|see{general relativity}}\index{general relativity},
elegantly embodied in his field equations.

% DOCX paragraph 106
Thus, we will learn to connect what we know and experience on our planet, where
gravitational fields are very weak, to what we observe in the universe, where enormously
stronger gravitational fields give life to marvellous phenomena such as black holes,
neutrons stars, and gravitational
waves\index{waves!gravitational (see gravitational wave)}. In this way, we will find
that our knowledge and understanding of gravity is strongly influenced by the way it
manifests itself on this planet. However, we will also realise that we have to abandon
such a view because it is wrong and because the constraints it sets on our imagination
are too tight.

% DOCX paragraph 107
All in all, the correct answer to the questions posed above will be clear: ``Gravity is
simply the manifestation of the curvature of spacetime''. At the moment, this statement
is still cryptic, I know. But, as mentioned in the introduction, a certain amount of
patience is needed to embark on the journey that awaits us. I can assure you that it
will all become much clearer by the end of Chapter 3.
